\documentclass[letterpaper,10pt]{article}

\usepackage[letterpaper,margin=0.48in]{geometry}
\usepackage[T1]{fontenc}
\usepackage[scaled=0.96]{helvet}
\usepackage{xcolor}
\usepackage{hyperref}
\usepackage{enumitem}
\usepackage[expansion=false]{microtype}

\definecolor{cobalt}{HTML}{1D4ED8}
\renewcommand{\familydefault}{\sfdefault}
\hypersetup{
  colorlinks=true,
  urlcolor=cobalt,
  linkcolor=cobalt
}

\pagestyle{empty}
\setlength{\parindent}{0pt}
\setlength{\parskip}{0pt}
\setlength{\emergencystretch}{2em}
\setlist[itemize]{
  leftmargin=1.2em,
  label=\textbullet,
  itemsep=0.8pt,
  topsep=1.5pt,
  parsep=0pt,
  partopsep=0pt
}

\newcommand{\sectiontitle}[1]{%
  \vspace{3.5pt}%
  {\color{cobalt}\large\bfseries #1}\par
  \vspace{1pt}\hrule height 0.35pt
  \vspace{2.5pt}%
}
\newcommand{\datedtitle}[2]{%
  \noindent\textbf{#1}\hfill{\small #2}\par
}
\newcommand{\role}[3]{%
  \datedtitle{#1}{#2}%
  \textit{#3}\par
}
\newcommand{\pubitem}[2]{%
  \item[\textbf{#1}] #2
}

\begin{document}

\begin{center}
  {\Huge\bfseries\color{cobalt} SHIVA SHRESTHA}\par
  \vspace{1pt}
  {\small Applied Machine Learning | Embedded Systems | Secure Edge AI}\par
  \vspace{3pt}
  {\scriptsize
    \href{mailto:sshres16@students.kennesaw.edu}{sshres16@students.kennesaw.edu}
    \textbar\ 
    \href{mailto:sh7vashrestha@gmail.com}{sh7vashrestha@gmail.com}
    \textbar\ 
    \href{https://linkedin.com/in/sh7va}{linkedin.com/in/sh7va}
    \textbar\ 
    \href{https://github.com/sh7vashrestha}{github.com/sh7vashrestha}
    \textbar\ 
    \href{https://shivashrestha7.com.np}{https://shivashrestha7.com.np}
    \textbar\ Marietta, Georgia, USA
  }\par
\end{center}
\vspace{-4pt}

\sectiontitle{SUMMARY}
Applied ML researcher specializing in secure on-device LLMs, multimodal medical AI, and embedded intelligent systems. Experienced in LLM fine-tuning, federated optimization, differential privacy, edge deployment, robotics, IoT, and full-stack embedded ML pipelines.

\sectiontitle{EDUCATION}
\datedtitle{PhD, Computer Science}{2025--Present}
Kennesaw State University, USA\\
Focus: Edge Security, Federated Learning, On-device LLMs\\
Advisor: Dr. Honghui Xu

\vspace{2pt}
\datedtitle{B.E., Electronics, Communication \& Information Eng.}{2019--2024}
IOE WRC, Tribhuvan University\\
82.32\% (Faculty Topper)

\sectiontitle{RESEARCH \& INDUSTRY EXPERIENCE}
\role{Graduate Research Assistant}{2025--Present}{Kennesaw State University}
\begin{itemize}
  \item UAV-MAID: Built an intrusion detection framework for drone swarms that judges each drone from what its honest neighbours overhear it relay, fusing a random forest, an attention-based relay network and a flight-relative outlier detector. On 567 simulated ns-3 flights it detected 80\% of attack types unseen in training, against 59\% for the best of ten baselines.
  \item KESTREL: Built an offline AI search-and-rescue drone on a Jetson Orin Nano and Pixhawk 6X. An on-device Gemma model plans the survey from a phone map or chat request, and YOLO, a thermal camera and CLIP find the person, with human approval gates before flight.
  \item DP-FedLoRA: Contributed to differentially private federated on-device LLM adaptation; recognized as ICDM 2025 Best Paper Runner-Up.
  \item KDLLM (First Author): Developed a knowledge-distillation-based framework for LLM copyright preservation, enabling compact student-model training while protecting proprietary teacher-model ownership.
  \item MedSpeak: Built a spoken medical QA pipeline integrating Whisper ASR, knowledge-graph grounding, and LLM-based semantic correction.
  \item ClinX: Designed a privacy-preserving multimodal medical VQA framework with OCR-based PHI masking, PP-GAN restoration, and LLaVA-Med evaluation.
  \item Edge-LLM IDS: Implemented syscall-based anomaly detection using n-gram features, Isolation Forest, and LLM-assisted explanation on Raspberry Pi.
  \item DamageScout (UAV Damage Assessment): Built an onboard computer-vision pipeline for real-time disaster damage assessment on Jetson Orin Nano, integrating GPS geo-tagging and Pixhawk telemetry. It groups repeated detections into building-level damage records.
\end{itemize}

\role{Electronics Engineer}{Apr 2024--Aug 2024}{Yarsa Tech, Nepal}
\begin{itemize}
  \item Built STM32 multilingual toy: Flash storage + DAC + amplifier audio system.
  \item Developed ESP8266 audio playback system with SPI Flash.
  \item Designed Smart Cloud Printer: Raspberry Pi + MQTT + CUPS for automated legacy printing.
  \item Automated industrial barcode printing (Python + Linux).
\end{itemize}

\sectiontitle{PUBLICATIONS}
\textbf{Journal Articles}
\begin{itemize}[leftmargin=2.2em,label={}]
  \pubitem{[J1]}{\textbf{S. Shrestha}, H. Xu, Z. Xie, D. Seo, Y. Joe, W. Kim, and Y. Li. KDLLM: Copyright-Preserving LLM based on Knowledge Distillation. Tsinghua Science and Technology, 2025. \href{https://www.sciopen.com/article/10.26599/TST.2025.9010130}{Paper Link}}
  \pubitem{[J2]}{H. Xu, K. Li, \textbf{S. Shrestha}, D. Zheng, W. Chen, Z. Li, and Z. Cai. A Survey: Towards Privacy and Security in Mobile Large Language Models. Big Data Mining and Analytics, Accepted/In Press.}
  \pubitem{[J3]}{P. Paudel, Samman Shrestha, \textbf{Shiva Shrestha}, Sudarshan Gurung, and Smita Adhikari. Automated Waste Sorting with Delta Arm and YOLOv8 Detection. IRO Journal, 2024. \href{https://doi.org/10.36548/jaicn.2024.3.004}{Paper Link}}
\end{itemize}

\textbf{Conference Papers}
\begin{itemize}[leftmargin=2.2em,label={}]
  \pubitem{[C1]}{H. Xu, \textbf{S. Shrestha}, W. Chen, Z. Li, and Z. Cai. DP-FedLoRA: Privacy-Enhanced Federated Fine-Tuning for On-Device Large Language Models. ICDM 2025, Best Paper Runner-Up. \href{https://icdm.zhonghuapu.com/Awards/Awards.shtml}{Award Link}}
  \pubitem{[C2]}{Y. Song, \textbf{S. Shrestha}, C. Lyu, E. Khatibi, P. Zhang, H. Xu, N. Dutt, and A. Rahmani. MedSpeak: A Knowledge Graph-Aided ASR Error Correction Framework for Spoken Medical QA. ICASSP 2026, Accepted. \href{https://ieeexplore.ieee.org/document/11464058}{Paper Link}}
  \pubitem{[C3]}{\textbf{Shiva Shrestha}, Samman Shrestha, Prateek Paudel, Sudarshan Gurung, Sulav Gairee, and Smita Adhikari. Deep Learning for Waste Management: Leveraging YOLO for Accurate Waste Classification. IOE Graduate Conference, 2024. \href{http://conference.ioe.edu.np/publications/ioegc15/IOEGC-15-004-A1-4-61.pdf}{Paper Link}}
\end{itemize}

\textbf{Submitted Manuscripts}
\begin{itemize}[leftmargin=2.2em,label={}]
  \pubitem{[S1]}{\textbf{S. Shrestha}, K. S. Sharif, A. Xiang, Z. Xie, J. Huang, and H. Xu. Thermo-FL: Thermal-Aware Robust Federated Fine-Tuning of Large Language Models for Edge AI. Submitted to NDSS.}
  \pubitem{[S2]}{\textbf{S. Shrestha}, Z. Xie, C. Zhao, L. Ma, Z. Cai, and H. Xu. Masking Is Not Enough: Generative Restoration for Multimodal De-Identification in Medical AI. Submitted to NeurIPS.}
\end{itemize}

\sectiontitle{SELECTED PROJECTS}
\begin{itemize}
  \item UAV-MAID (Drone Swarm Intrusion Detection): Built a drone-swarm intrusion detector using honest-neighbour relay observations, a random forest, an attention-based relay network, and a flight-relative outlier detector. It detected 80\% of unseen attack types on 567 simulated ns-3 flights, against 59\% for the best of ten baselines.
  \item KESTREL (Offline Search-and-Rescue Drone, AGI Hackathon 2026): Built an offline drone on Jetson Orin Nano and Pixhawk 6X. An on-device Gemma model plans surveys from a phone map or chat request, while YOLO, a thermal camera, and CLIP find the person, with human approval gates before flight.
  \item MedSpeak (LLM + ASR + KG): Developed a spoken medical QA framework using Whisper ASR, knowledge-graph grounding, and LLM-based semantic error correction.
  \item ClinX: Designed a privacy-aware multimodal medical AI pipeline using PP-GAN-based image restoration and LLaVA-Med-based utility evaluation.
  \item Thermo-FL: Built a thermal-aware federated fine-tuning framework for edge LLMs with sparse LoRA communication and robust aggregation.
  \item Edge-LLM Intrusion Detection: Implemented syscall anomaly detection using n-gram feature extraction, Isolation Forest, and LLM-assisted analysis.
  \item DamageScout: Developed an edge-deployed visual damage assessment system for disaster mapping with onboard inference and GPS-based geo-tagging.
  \item Delta Arm Waste Classifier: Built a YOLOv8-based robotic waste-sorting system using a delta arm, achieving 0.98 mAP@50.
  \item Ambulance Triage (Sajilo): Developed an ambulance dispatch and triage application with SVM- and Random Forest-based severity prediction.
  \item LoRa Weather System: Created a long-range environmental telemetry system with LoRa communication and ML-based forecasting.
  \item Smart Cloud Printer: Built an IoT-enabled cloud printing system using Raspberry Pi, MQTT, and CUPS.
  \item STM32 Storytelling Toy: Developed an embedded multilingual audio playback system using STM32 Flash storage.
  \item Smart Access System: Designed an ESP32-based RFID access control system with integrated smoke detection and automation.
\end{itemize}

\sectiontitle{TECHNICAL SKILLS}
\begin{itemize}
  \item Languages: Python, C/C++, JavaScript
  \item Machine Learning: PyTorch, TensorFlow, Scikit-learn, PEFT/LoRA, Whisper, LLaVA, CLIP, XGBoost
  \item Edge and Embedded Systems: ArduPilot/MAVLink, NVIDIA Jetson, Raspberry Pi, STM32, ESP32/8266, Pixhawk, Embedded Linux
  \item IoT and Systems: MQTT, CUPS, syscall tracing, shell scripting, Docker, Linux, ns-3
  \item Web and Databases: React, Node.js, Tailwind, MongoDB, MySQL, Firebase
  \item Hardware and Development Tools: Git, KiCad, Keil, VS Code
\end{itemize}

\sectiontitle{HONORS \& AWARDS}
\begin{itemize}
  \item Best Paper Runner-Up, IEEE ICDM 2025 (DP-FedLoRA)
  \item Faculty Topper, BEI, IOE WRC (2024)
  \item Winner: NSU Demo (2024), E-GEN Expo (2024), MechTRIX (2023)
  \item Runner-Up: ICES Datathon (2024)
  \item District Topper, NEB +2 Science (2017)
\end{itemize}

\sectiontitle{LEADERSHIP \& ACTIVITIES}
\begin{itemize}
  \item Senior Member, Robotics Club, IOE WRC (2023--2024)
  \item Design Transformer, Robotics Booster, IOE WRC (2022--2023)
  \item Member, Robotics Club, IOE WRC (2021--2022)
  \item Coordinator, Battle for Speed Robotics (2022)
\end{itemize}

\end{document}
